\chapter{Tinjauan Pustaka}

\section{Penelitian Terkait}

Penelitian mengenai arsitektur multitenant pada platform berbasis web dan \textit{Software as a
Service} telah banyak dilakukan seiring meningkatnya kebutuhan sistem yang melayani banyak
organisasi dalam satu aplikasi terpusat. Fokus penelitian terdahulu umumnya mencakup perbandingan
model \textit{multitenancy}, mekanisme isolasi data, serta implikasi pemilihan arsitektur terhadap
skalabilitas dan efisiensi infrastruktur. Bagian ini merangkum beberapa penelitian relevan yang
menjadi landasan dalam perancangan arsitektur multitenant pada sistem manajemen aset, inventori,
dan rental.

Mangwani et al. (2023) \cite{ref5} melakukan evaluasi terhadap penerapan arsitektur multitenant
pada aplikasi SaaS dengan membandingkan pendekatan monolitik dan \textit{microservices}. Penelitian
tersebut menunjukkan bahwa arsitektur multitenant mampu meningkatkan efisiensi penggunaan sumber
daya dibandingkan pendekatan \textit{single tenant}, terutama ketika jumlah tenant meningkat.
Namun, studi ini lebih menekankan aspek \textit{deployment} dan performa layanan secara umum,
tanpa membahas secara spesifik implikasi pemilihan model penyimpanan data seperti \textit{shared
schema} terhadap isolasi data dan pengelolaan tenant.

Penelitian oleh Mangwani dan Tokekar (2022) \cite{ref9} menganalisis perilaku sistem SaaS dengan
pendekatan \textit{tenant shared} dan \textit{tenant specific}. Hasil penelitian menunjukkan bahwa
model \textit{tenant shared} lebih efisien dari sisi infrastruktur, tetapi memiliki risiko yang
lebih tinggi terhadap kebocoran data apabila mekanisme isolasi tidak dirancang dengan baik. Studi
ini menegaskan pentingnya pengelolaan konteks tenant dan pengendalian akses pada lapisan aplikasi,
namun tidak membahas bagaimana mekanisme tersebut dirancang secara terintegrasi dalam konteks
sistem manajemen aset dan inventori.

Levchenko dan Taratukhin (2021) \cite{ref7} mengusulkan metode adaptasi arsitektur multitenant
berbasis proses bisnis untuk mendukung fleksibilitas penggunaan sistem oleh berbagai organisasi.
Penelitian ini menyoroti pentingnya desain proses \textit{onboarding} dan pengelolaan tenant yang
terstandarisasi agar sistem dapat digunakan oleh tenant dengan karakteristik operasional yang
berbeda. Meskipun relevan, penelitian ini lebih berfokus pada adaptasi proses bisnis, bukan pada
perancangan arsitektur penyimpanan data berbasis \textit{shared database} dan \textit{shared
schema}.

Sementara itu, Pahl dan Lee (2020) \cite{ref10} membahas pola arsitektur \textit{multitenancy}
pada sistem SaaS dan mengidentifikasi \textit{trade-off} antara model \textit{separate database},
\textit{separated schema}, dan \textit{shared schema}. Penelitian ini memberikan landasan teoritis
mengenai implikasi masing-masing model terhadap isolasi data dan efisiensi infrastruktur. Namun,
kajian tersebut bersifat konseptual dan belum mengaitkan secara langsung penerapan \textit{shared
schema} dengan konteks sistem manajemen aset, inventori, dan rental pada platform multi organisasi.

Berdasarkan penelitian terdahulu tersebut, dapat disimpulkan bahwa sebagian besar studi berfokus
pada evaluasi umum arsitektur multitenant, perbandingan model \textit{deployment}, atau adaptasi
proses bisnis. Masih terbatas penelitian yang secara spesifik membahas perancangan arsitektur
multitenant berbasis \textit{shared database} dan \textit{shared schema} dengan fokus pada
efisiensi infrastruktur serta mekanisme isolasi data dalam sistem manajemen aset, inventori, dan
rental. Oleh karena itu, penelitian ini diarahkan untuk mengisi celah tersebut dengan merancang
arsitektur multitenant yang secara eksplisit mengintegrasikan pemilihan model penyimpanan data,
mekanisme isolasi tenant, dan kebutuhan pengelolaan sistem multi organisasi dalam satu platform
terpusat.

\section{Sistem Informasi Manajemen Aset dan Inventori}

Sistem informasi manajemen aset dan inventori merupakan perangkat lunak yang digunakan untuk
melakukan pencatatan, pemantauan, pengendalian, dan pelaporan aset serta barang milik organisasi
secara terstruktur. Penerapan sistem informasi pada pengelolaan aset terbukti mampu meningkatkan
akurasi data, transparansi operasional, serta efisiensi pemanfaatan sumber daya dibandingkan
pendekatan manual atau pencatatan terpisah \cite{ref1,ref3}. Dalam sistem yang terintegrasi,
informasi mengenai kepemilikan aset, kondisi barang, serta riwayat penggunaan dapat diakses
secara konsisten dan \textit{real time}.

Pada platform yang digunakan oleh banyak organisasi secara bersamaan, sistem manajemen aset dan
inventori menghadapi kompleksitas tambahan dibandingkan sistem satu organisasi. Setiap organisasi
memiliki struktur data, pengguna, dan alur operasional yang berbeda, sehingga pengelolaan data
aset tidak hanya berkaitan dengan fungsi pencatatan, tetapi juga dengan mekanisme pemisahan data
dan pengendalian akses. Tanpa dukungan arsitektur yang tepat, pengelolaan aset pada platform multi
organisasi berpotensi menimbulkan duplikasi data, inkonsistensi informasi, serta kesulitan dalam
pengelolaan sistem secara terpusat.

Oleh karena itu, integrasi sistem manajemen aset, inventori, dan layanan rental pada platform
multi organisasi tidak dapat dipisahkan dari perancangan arsitektur perangkat lunak yang mampu
mendukung pengelolaan data lintas organisasi secara efisien dan terisolasi. Aspek ini menjadikan
sistem manajemen aset dan inventori relevan untuk dikaji tidak hanya dari sisi fungsional, tetapi
juga dari perspektif arsitektur sistem.

\section{Konsep \textit{Cloud Computing} dan SaaS}

\textit{Cloud computing} menyediakan infrastruktur komputasi yang elastis, terdistribusi, dan
dapat diakses secara \textit{on demand}, sehingga memungkinkan pengembangan dan pengelolaan
aplikasi berskala besar dengan efisiensi sumber daya yang lebih baik \cite{ref2}. Dalam konteks
\textit{Software as a Service}, aplikasi disediakan sebagai layanan yang dapat diakses melalui
jaringan tanpa memerlukan instalasi lokal, dengan model penggunaan yang mendukung banyak pengguna
dan organisasi secara bersamaan.

Arsitektur SaaS memungkinkan satu aplikasi digunakan oleh banyak organisasi melalui satu sistem
terpusat, sehingga mengurangi biaya operasional, menyederhanakan proses pemeliharaan, dan
mempermudah pembaruan sistem \cite{ref5}. Namun, implementasi SaaS tidak selalu identik dengan
arsitektur multitenant. Banyak sistem SaaS masih dibangun menggunakan pendekatan non multitenant
atau \textit{single tenant}, di mana setiap organisasi memiliki konfigurasi atau \textit{instance}
sistem tersendiri. Pendekatan tersebut berpotensi menimbulkan keterbatasan skalabilitas dan
efisiensi ketika jumlah organisasi pengguna meningkat.

Dalam konteks platform layanan multi organisasi berskala nasional, pemilihan arsitektur SaaS yang
tepat menjadi krusial. Arsitektur yang tidak dirancang untuk penggunaan bersama oleh banyak tenant
akan meningkatkan kompleksitas pengelolaan sistem dan mengurangi manfaat utama dari model SaaS
itu sendiri.

\section{Arsitektur Multitenant}

Arsitektur multitenant merupakan pendekatan perancangan aplikasi di mana satu \textit{instance}
aplikasi digunakan oleh banyak tenant secara bersamaan, dengan tetap menjaga isolasi data,
konfigurasi, dan proses operasional antar tenant \cite{ref4}. Tenant dalam arsitektur multitenant
didefinisikan sebagai entitas organisasi independen yang memiliki data, pengguna, dan alur kerja
tersendiri dalam satu sistem bersama.

Berbagai penelitian mengidentifikasi beberapa pola utama dalam penerapan \textit{multitenancy},
yaitu \textit{shared schema}, \textit{separated schema}, dan \textit{separate database}, yang
masing-masing memiliki karakteristik dan \textit{trade-off} pada aspek isolasi data, biaya
infrastruktur, dan skalabilitas sistem \cite{ref6}. Model \textit{separate database} menawarkan
isolasi data yang tinggi, namun membutuhkan sumber daya dan biaya pengelolaan yang besar.
Sebaliknya, model \textit{shared schema} menawarkan efisiensi infrastruktur dan kemudahan
pengelolaan, tetapi memerlukan mekanisme isolasi data yang lebih ketat pada lapisan aplikasi.

Model \textit{shared database} dan \textit{shared schema} banyak direkomendasikan untuk platform
SaaS tahap awal dan sistem berskala besar karena mampu mendukung pertumbuhan jumlah tenant dengan
biaya infrastruktur yang relatif rendah serta menyederhanakan proses pemeliharaan sistem
\cite{ref9,ref10}. Pemilihan model ini selaras dengan kebutuhan efisiensi infrastruktur dan
pengelolaan data banyak organisasi dalam satu platform terpusat, sebagaimana menjadi fokus
penelitian ini.

\section{Isolasi Data pada Sistem Multitenant}

Isolasi data merupakan aspek kritis dalam arsitektur multitenant untuk memastikan bahwa setiap
tenant hanya dapat mengakses data yang menjadi miliknya. Isolasi data dapat diterapkan pada
beberapa lapisan sistem, yaitu lapisan basis data, lapisan aplikasi, dan lapisan infrastruktur
\cite{ref5,ref8,ref9}.

Pada lapisan basis data, isolasi dapat dilakukan melalui penggunaan atribut identitas tenant
seperti \texttt{tenant\_id}, pemisahan skema, atau pemisahan basis data. Pada lapisan aplikasi,
isolasi diwujudkan melalui mekanisme \textit{tenant aware middleware}, pengelolaan konteks tenant,
serta pembatasan akses berbasis peran. Sementara itu, pada lapisan infrastruktur, isolasi
dilakukan melalui segmentasi sumber daya dan pengendalian beban kerja per tenant.

Penelitian empiris menunjukkan bahwa pendekatan \textit{tenant shared} lebih efisien secara biaya
dan sumber daya dibandingkan pendekatan \textit{tenant specific}, namun memerlukan desain isolasi
logis yang ketat untuk menjaga keamanan dan konsistensi data \cite{ref9}. Oleh karena itu, dalam
konteks arsitektur multitenant berbasis \textit{shared schema}, fokus isolasi data lebih diarahkan
pada pengelolaan konteks tenant dan \textit{tenant scoping} pada lapisan aplikasi, yang menjadi
salah satu fokus utama penelitian ini.

\section{\textit{Tenant Provisioning}}

\textit{Tenant provisioning} merupakan proses pembuatan dan inisialisasi tenant baru dalam sistem
multitenant, yang meliputi konfigurasi awal, inisialisasi metadata, pembentukan entitas organisasi,
serta penetapan peran pengguna awal \cite{ref7}. Dalam sistem non multitenant, proses ini umumnya
dilakukan secara manual dan terpisah untuk setiap organisasi, sehingga meningkatkan kompleksitas
operasional dan risiko ketidakkonsistenan konfigurasi.

Pada arsitektur multitenant, \textit{provisioning} tenant yang terstandarisasi dan otomatis
menjadi bagian penting dari efisiensi operasional sistem. Otomatisasi \textit{provisioning} tenant
memungkinkan platform untuk menambahkan organisasi baru tanpa memerlukan konfigurasi manual
tambahan, sekaligus memastikan konsistensi struktur data dan kebijakan akses antar tenant. Dengan
demikian, mekanisme \textit{provisioning} tenant tidak hanya bersifat administratif, tetapi juga
berperan sebagai komponen arsitektural yang mendukung skalabilitas dan efisiensi sistem secara
keseluruhan.
